\documentclass[11pt]{article}
\usepackage[margin=1in]{geometry}
\usepackage{amsmath,amssymb}
\usepackage{booktabs}
\usepackage{graphicx}
\usepackage[colorlinks=true,linkcolor=blue,urlcolor=blue]{hyperref}

\newcommand{\arcsecs}{\ensuremath{''}}

\title{Proposed Use of Engineering Time:\\
Rate-Matched Imaging of the Geosynchronous Belt}
\author{C.~Stubbs, on behalf of the GEO luminosity-function team}
\date{\today}

\begin{document}
\maketitle

\section{Project goal}

The science objective is a \textbf{deep luminosity function of objects in and around the
geosynchronous (GEO) belt} --- active satellites, dead payloads, and debris --- extending
well below the completeness limit of current catalogs. Depth requires long effective
integrations; long integrations on GEO objects require tracking that \emph{freezes the
targets, not the stars}.

\section{Engineering-time objectives}

\begin{enumerate}
  \item Obtain \textbf{representative frames containing known artificial satellites} across
        a range of orbital inclinations, to exercise and hone the data-reduction pipeline
        (streaked-star astrometry, target photometry, co-addition).
  \item \textbf{Validate and refine the observing strategy} described below --- in
        particular the inclination-dependent rate-lock budget of
        \S\ref{sec:ratelock}, which fixes the exposure time, sequence length, and the
        readout overhead we can tolerate.
\end{enumerate}

\section{Observing strategy}

We will \textbf{rate-match in right ascension as a function of inclination}: the telescope
tracks at the apparent RA rate of a geosynchronous object of the chosen inclination class,
freezing satellites and debris of that class in the image. Stars streak; we have
previously demonstrated the ability to extract \textbf{reliable centroids and precision
astrometry from stellar streaks}, so the astrometric solution is not degraded.

Observations are scheduled at each target's \textbf{southern declination peak} (the
culmination of its daily analemma arc), for three reasons:
(i) the declination rate passes through zero there, so a pure-RA rate match is exact at
mid-sequence; (ii) the RA rate is itself at an extremum there (\S\ref{sec:ratelock}),
so it holds constant to high accuracy; and (iii) the southern turning point maximizes
angular separation from the crowded zero-inclination belt.

At each tracking rate we will take a \textbf{succession of short exposures for
co-addition, totalling $T_{\rm int} = 200$~s} of integration. Field centers are offset
north so the target sits $75\%$ of the way toward the \emph{southern} edge of the field
of view at culmination; the object then climbs back north through the field, maximizing
dwell time. Nightly target lists (LMST-ordered peak times, pointings in RA/Dec/HA/Alt/Az,
and tracking rates with signs defined) are generated by the accompanying
\texttt{geo\_peaks.py} planning tool from same-day TLEs.

\paragraph{Sign conventions.} RA rate $>0$ means RA increasing (eastward motion).
The sidereal reference rate is $\dot\alpha_{\rm sid}=15.0411\arcsecs$/s (the LMST rate,
i.e.\ the RA rate of an ideal zero-inclination geostationary point). We quote object
rates and the offset $\Delta \equiv \dot\alpha - \dot\alpha_{\rm sid}$; $\Delta>0$ means
the object drifts \emph{east} of a geostationary point. At the southern peak an inclined
object runs eastward-fast: $\Delta \simeq f\,\dot\alpha_{\rm sid}\,(\sec i - 1)$, where
$f\simeq1.17$ is the measured topocentric amplification from CTIO (Table~\ref{tab:meas}).

\paragraph{Night thirds and illumination.} The night (Sun below $-12^{\circ}$) is split
into thirds, and within each third targets are ranked by \textbf{solar phase angle}
(Sun--object--observer; smaller = better lit). Because the anti-solar point rises in the
east at dusk, crosses the meridian at midnight, and sets in the west at dawn, this ranking
automatically points the first third \textbf{east}, the middle third near the
\textbf{meridian}, and the last third \textbf{west}. The planner also flags objects that
sit inside \textbf{Earth's shadow} at their peak (a cylindrical-umbra test); eclipsed
targets are excluded from the plan.

\section{Passband selection}
\label{sec:band}

\textbf{Rule: $g$ band when the Moon is down; $r$ band in bright time.}
Moonlight is scattered sunlight, and Rayleigh scattering scales as $\lambda^{-4}$: the
lunar sky brightens $g$ by $\gtrsim 2$~mag between dark and bright time, but $r$ by only
$\sim$1~mag. In dark time the CTIO sky is faintest in $g$
($\mu_g \approx 22.0$ vs.\ $\mu_r \approx 21.0$~mag\,arcsec$^{-2}$), and $g$ sits near
peak system throughput, so $g$ delivers the deepest limiting magnitude. Once the Moon is
up, the $g$ sky floods and $r$ becomes the depth optimum. A secondary effect points the
same way: satellite reflectance spectra (solar panels, MLI, aged coatings) are typically
\emph{reddened} relative to solar, so the targets themselves yield a small additional
advantage in $r$ that costs little in dark time and helps in bright time.

Two operational cautions. First, the crossover is governed by \textbf{Moon--target
separation}, not Moon phase alone: a gibbous Moon $40^{\circ}$ from the field damages $g$
nearly as much as a full Moon $90^{\circ}$ away. Second, the illumination scheduling of
\S3 prefers \emph{anti-solar} pointings --- and at full Moon the anti-solar point is
\emph{at the Moon}. In bright time we therefore back off from the lowest-phase-angle
pointing as needed to hold $\gtrsim 30$--$40^{\circ}$ of lunar separation, accept a
modestly larger phase angle, and observe in $r$. The engineering run will measure the
actual $g$/$r$ depth crossover directly (\S6, item 5).

\section{Rate-lock analysis: how long does the match hold?}
\label{sec:ratelock}

For a circular geosynchronous orbit of inclination $i$ observed at its southern peak
(time $t$ measured from culmination, $n = 2\pi/86164\,{\rm s}$):

\medskip
\noindent\textbf{RA is not the constraint.} The apparent RA rate is at an extremum at
the peak; its residual after matching at $t=0$ grows only cubically,
\begin{equation}
\delta\alpha(t) \;=\; \frac{n^{3}\sin^{2}i}{3\cos^{3}i}\,t^{3}
\qquad\Longrightarrow\qquad
\delta\alpha(\pm300\,{\rm s}) \;<\; 0.05\arcsecs
\;\;\text{even at } i=14^{\circ}.
\end{equation}
A single commanded RA rate therefore holds for the whole sequence.

\medskip
\noindent\textbf{Declination drift sets the budget.} The declination rate is zero at the
peak but accelerates away from it:
\begin{equation}
\ddot\delta \;=\; k(i) \;=\; f\, i\, n^{2}
\qquad\text{(topocentric factor } f \simeq 1.17 \text{ from CTIO, measured)} .
\end{equation}
We verified $k$ numerically with current TLEs propagated topocentrically from CTIO
(Table~\ref{tab:meas}); the measured values match $f\,i\,n^{2}$ to a few percent.
An exposure of length $\tau$ taken at time $t$ from the peak smears the target in
declination by $k\,|t|\,\tau$. Requiring per-exposure smear below a budget
$\epsilon$ confines the sequence to the window
\begin{equation}
|t| \;\le\; \frac{\epsilon}{k(i)\,\tau}
\qquad\Longrightarrow\qquad
W \;=\; \frac{2\epsilon}{k(i)\,\tau}
\quad\text{(total usable wall time, centered on the peak).}
\end{equation}
With $N=T_{\rm int}/\tau$ exposures the sequence needs wall time $N(\tau+R)$ for
readout time $R$, so the \textbf{maximum tolerable readout} is
\begin{equation}
R_{\max} \;=\; \frac{W}{N} - \tau \;=\; \frac{2\epsilon}{k(i)\,T_{\rm int}}\,\tau - \tau .
\label{eq:rmax}
\end{equation}

\begin{table}[ht]
\centering
\caption{Measured declination acceleration $k$ at the southern peak, from current TLEs
propagated topocentrically from CTIO (2026-08-14). The implied topocentric factor
$f = k/(i\,n^{2})$ is constant at $\simeq 1.17$.}
\label{tab:meas}
\begin{tabular}{lccc}
\toprule
Object & $i$ (deg) & $k$ ($\arcsecs\,{\rm s}^{-2}$) & implied $f$ \\
\midrule
ANIK F2        &  3.43 & $7.5\times10^{-5}$ & 1.15 \\
GALAXY 3C      &  3.88 & $8.7\times10^{-5}$ & 1.17 \\
DIRECTV 5      &  5.02 & $1.10\times10^{-4}$ & 1.14 \\
INTELSAT 902   &  6.26 & $1.40\times10^{-4}$ & 1.17 \\
INTELSAT 904   &  7.06 & $1.54\times10^{-4}$ & 1.14 \\
NIMIQ 2        &  9.58 & $2.09\times10^{-4}$ & 1.14 \\
TDRS 3         & 12.56 & $2.81\times10^{-4}$ & 1.17 \\
USA 115        & 14.09 & $3.18\times10^{-4}$ & 1.18 \\
TDRS 6         & 14.18 & $3.19\times10^{-4}$ & 1.17 \\
\bottomrule
\end{tabular}
\end{table}

\begin{table}[ht]
\centering
\caption{Rate-lock budget for $T_{\rm int}=200$~s: usable window $W$ and maximum
tolerable readout $R_{\max}$ (Eq.~\ref{eq:rmax}), for per-exposure declination smear
budget $\epsilon$ and exposure time $\tau$. Negative $R_{\max}$ means 200~s of
integration \emph{cannot} be completed in a single peak visit at that configuration.}
\label{tab:budget}
\begin{tabular}{cc|cc|cc|cc}
\toprule
& & \multicolumn{2}{c|}{$\epsilon=0.5\arcsecs$, $\tau=20$ s}
  & \multicolumn{2}{c|}{$\epsilon=0.5\arcsecs$, $\tau=10$ s}
  & \multicolumn{2}{c}{$\epsilon=1.0\arcsecs$, $\tau=20$ s} \\
$i$ (deg) & $k$ ($\arcsecs/{\rm s}^2$) & $W$ (s) & $R_{\max}$ (s) & $W$ (s) & $R_{\max}$ (s) & $W$ (s) & $R_{\max}$ (s) \\
\midrule
 3 & $6.7\times10^{-5}$  &  744 & $+54$ & 1488 & $+64$ & 1488 & $+129$ \\
 5 & $1.1\times10^{-4}$  &  446 & $+25$ &  893 & $+35$ &  893 & $+69$  \\
 8 & $1.8\times10^{-4}$  &  279 & $+8$  &  558 & $+18$ &  558 & $+36$  \\
10 & $2.2\times10^{-4}$  &  223 & $+2$  &  446 & $+12$ &  446 & $+25$  \\
12 & $2.7\times10^{-4}$  &  186 & $-1$  &  372 & $+9$  &  372 & $+17$  \\
14 & $3.1\times10^{-4}$  &  159 & $-4$  &  319 & $+6$  &  319 & $+12$  \\
\bottomrule
\end{tabular}
\end{table}

\paragraph{Conclusions from the budget.}
\begin{itemize}
  \item \textbf{Low inclination ($i \lesssim 5^{\circ}$):} 20-s exposures with generous
        readout ($R\lesssim25$~s) comfortably reach 200~s in one visit. Readout losses
        are not a driver here.
  \item \textbf{Mid inclination ($i \sim 8$--$10^{\circ}$):} 20-s exposures work only if
        readout is fast ($R\lesssim 2$--$8$~s). With a DECam-class readout
        ($R \approx 20$~s), \emph{shorter exposures are required}: at $\tau=10$~s the
        budget allows $R_{\max}\approx12$--$18$~s --- marginal. This is precisely the
        regime the engineering run must probe.
  \item \textbf{High inclination ($i \gtrsim 12^{\circ}$):} a single peak visit cannot
        deliver 200~s at $\epsilon=0.5\arcsecs$, $\tau=20$~s at all ($W<200$~s). Options,
        in order of preference: (i)~\textbf{command a per-frame declination rate}
        ($\dot\delta = k\,t$, a linear ramp): the residual per exposure drops to
        $\tfrac12 k\tau^{2}\lesssim0.06\arcsecs$ and the window is then limited only by
        the cubic RA term, i.e.\ many tens of minutes; (ii)~shorten exposures to
        $\tau=10$~s; (iii)~relax the smear budget to $1\arcsecs$; (iv)~split the sequence
        across two consecutive peak passages and co-add across visits.
  \item Because the drift is \emph{symmetric} about the peak, sequences must be
        \textbf{centered on the peak time} (start $W/2$ before culmination) --- the
        LMST-ordered peak times from \texttt{geo\_peaks.py} are the sequence midpoints.
  \item If co-addition registers each frame on the measured target position
        (shift-and-add), the accumulated declination \emph{offset} across the sequence
        (up to $\sim\!\tfrac12 k (W/2)^2 \approx 3$--$13\arcsecs$) is removed and only the
        per-exposure smear above matters; if instead fixed-WCS stacking is used to
        preserve star-streak astrometry, the offset itself must fit the smear budget and
        the windows above shrink by $\sqrt{2\epsilon/(kW\tau)}$. Testing both reduction
        paths is an explicit engineering goal.
\end{itemize}

\section{Proposed engineering sequences}

\begin{enumerate}
  \item \textbf{Rate ladder} (one visit each): known satellites at
        $i \approx \{0, 3.5, 5, 7, 9.6, 12.6, 14.2\}^{\circ}$ from the nightly
        \texttt{geo\_peaks} list (e.g.\ EchoStar~10, ANIK~F2, DIRECTV~5, INTELSAT~904,
        NIMIQ~2, TDRS~3, TDRS~6), each with a 200-s co-add sequence centered on its
        southern peak, at the matched RA rate and $\dot\delta=0$. Verifies the
        rate-lock windows of Table~\ref{tab:budget} directly against trailed-image
        metrology.
  \item \textbf{Declination-ramp test} at $i\approx12$--$14^{\circ}$: repeat with the
        per-frame $\dot\delta$ ramp, to qualify option~(i) above and unlock the
        high-inclination population for the science survey.
  \item \textbf{Exposure-time pair} at $i\approx10^{\circ}$: $\tau=10$~s vs $\tau=20$~s
        sequences of equal $T_{\rm int}$, to measure the real readout overhead and the
        smear--depth tradeoff in the pipeline.
  \item \textbf{Pipeline fields}: two zero-rate-offset fields on the crowded
        zero-inclination belt (streaked stars, frozen belt objects) for the
        luminosity-function reduction chain end-to-end.
  \item \textbf{$g$/$r$ depth crossover} (\S\ref{sec:band}): one field observed in both
        $g$ and $r$, in dark time and again with the Moon up at two lunar separations
        ($\sim40^{\circ}$ and $\sim90^{\circ}$), with equal $T_{\rm int}$ --- fixing
        empirically where the band switch belongs and quantifying the reddened-target
        advantage in $r$.
\end{enumerate}

Total request: the sequences above fit comfortably in \textbf{two half-nights} of
engineering time, including calibrations, with peak times spread across the LMST range
(the nightly list spans LMST $\approx 15$--$04$~h for the inclined fleet visible from
CTIO).

\section*{Tooling}

\texttt{geo\_peaks.py} (this directory) generates the nightly list from same-day TLEs
(CelesTrak or IAU CPS SatChecker): LMST-ordered southern-peak times, object and pointing
RA/Dec, hour angle, alt/az, RA rates and offsets from sidereal with the sign conventions
above, zenith angles, and sub-satellite longitudes, plus declination-track,
sky-position, and rate-offset plots. All numbers in \S\ref{sec:ratelock} were computed
with the same propagation code used for the nightly lists.

\end{document}
